\section{A common density for the two good regions}
\label{sec:replica_good_pair_marginal}

Let $Q,Y,V,C,R$ be finite coordinate sets and let $k\geq0$.
For $B\subseteq\{0,\ldots,k-1\}$, put $G=B^c$ and $m=|G|$.
Choose one enumeration of $G$ for all four good regions, and one enumeration
of $B$ for both bad auxiliary regions. For arbitrary
$\Omega\in\mathbb C^Q\otimes\mathbb C^Y\otimes\mathbb C^V$ and
$u\in(QYV)^{\otimes k}\otimes C^{\otimes k}\otimes R^{\otimes k}$,
put $w=(R_B\otimes I)u$. Here $R_B$ is the excitation operator determined
by $\Omega$; it is an orthogonal projection when $\Omega$ is unit.
The following coordinate identities trace every middle physical coordinate and identify
the two reduced matrices used
in \cite[Section~7]{OpenAI2026AreaLaw}, equation
(comparator:merge-moments). No permutation symmetry is required.

\begin{definition}[Common good-copy density]
\label{def:replica_good_pair_marginal}
\lean{Matrix.replicaGoodPairMarginal}
\leanok
\uses{def:good_bad_copies_equiv}
Regroup the same component $w$ into retained and discarded factors:
\begin{align}
 \mathcal H_{\mathrm{ret}}
   &=(Q^{\otimes m}\otimes C^{\otimes m})
      \otimes(V^{\otimes m}\otimes R^{\otimes m}),\\
 \mathcal H_{\mathrm{tr}}
   &=Y^{\otimes m}\otimes(QYV)^{\otimes|B|}\otimes C^{\otimes|B|}\otimes R^{\otimes|B|}.
\end{align}
The common density $\Gamma_B(u)$ is
$\operatorname{Tr}_{\mathcal H_{\mathrm{tr}}}|w\rangle\langle w|$.
It retains the four exterior good regions and is unnormalized. The discarded
factors include the good middle copies and the middle coordinates within the
bad physical copies, so every middle physical coordinate is traced.
\end{definition}

\begin{theorem}[First actual marginal]
\label{thm:replica_good_pair_marginal_right}
\lean{Matrix.partialTraceRight_replicaGoodPairMarginal}
\leanok
\uses{def:replica_good_pair_marginal,def:replica_good_regional_auxiliary_marginal}
The first marginal of the common density is the actual regional auxiliary
marginal of the same component:
\begin{align}
 \operatorname{Tr}_{V_G R_G}\Gamma_B(u)=\rho_{Q_G C_G}(w).
\end{align}
\end{theorem}
\begin{proof}
\leanok
Split the full $R$ configurations into their good and bad coordinates.
The two successive partial traces then sum over exactly the same coordinates
as the definition of $\rho_{Q_G C_G}(w)$. A finite-sum reordering proves the identity.
\end{proof}

\begin{theorem}[Second actual marginal]
\label{thm:replica_good_pair_marginal_left}
\lean{Matrix.partialTraceLeft_replicaGoodPairMarginal}
\leanok
\uses{def:replica_good_pair_marginal,def:replica_good_regional_auxiliary_marginal}
Exchange the exterior physical regions and the auxiliary regions simultaneously,
preserving the independent middle region:
\begin{align}
 \tau(q,(y,v))&=(v,(y,q)),\\
 \Omega'(v,(y,q))&=\Omega(q,(y,v)),\\
 u'(x,(r,c))&=u\bigl((i\mapsto\tau^{-1}(x_i)),(c,r)\bigr).
\end{align}
Let $w'=(R'_B\otimes I)u'$ be the actual component determined by $\Omega'$.
Then
\begin{align}
 \operatorname{Tr}_{Q_G C_G}\Gamma_B(u)=\rho_{V_G R_G}(w').
\end{align}
The vector $w'$ is the coordinate exchange of $w$, rather than a separately
chosen excitation component.
\end{theorem}
\begin{proof}
\leanok
The physical exchange preserves the middle coordinate and carries each factor of the excitation
operator to the corresponding factor for the exchanged vector. Tensoring
and applying it to the exchanged original vector therefore gives exactly
the exchanged component. The physical good/bad split respects this exchange.
Partial trace is invariant under the exchange of discarded coordinates,
so the common density is exchanged only on its retained factors. Apply
the first marginal identity to this exchanged density.
\end{proof}

\begin{theorem}[Actual component mass]
\label{thm:replica_good_pair_marginal_trace}
\lean{Matrix.trace_replicaGoodPairMarginal}
\leanok
\uses{def:replica_good_pair_marginal}
The common density has total mass
\begin{align}
 \operatorname{Tr}\Gamma_B(u)=\|w\|^2.
\end{align}
The trace is a complex number; the right-hand side is the real squared norm
included into $\mathbb C$. The identity includes zero components, zero copies
and empty good sets.
\end{theorem}
\begin{proof}
\leanok
The retained and discarded coordinates together form an equivalence with
the original coordinates. Reindexing by an equivalence and taking a partial
trace preserve the full trace. Finally,
$\operatorname{Tr}|w\rangle\langle w|=\langle w,w\rangle=\|w\|^2$.
\end{proof}
